% ============================================================
%  ChemoInteract Documentation - Part 1
%  Title: System Overview, Architecture, & Datasets
%  Compile with: pdflatex ChemoInteract_Part1_Overview_and_Architecture.tex
% ============================================================

\documentclass[12pt,a4paper,oneside]{article}

\usepackage[T1]{fontenc}
\usepackage[utf8]{inputenc}
\usepackage{lmodern}
\usepackage[margin=2.5cm]{geometry}
\usepackage{hyperref}
\usepackage{xcolor}
\usepackage{listings}
\usepackage{longtable}
\usepackage{booktabs}
\usepackage{array}
\usepackage{graphicx}
\usepackage{amsmath}
\usepackage{amssymb}
\usepackage{tikz}
\usepackage{tcolorbox}
\usepackage{fancyhdr}
\usepackage{titlesec}
\usepackage{enumitem}
\usepackage{microtype}
\usepackage{setspace}

\tcbuselibrary{breakable,skins}

\hypersetup{
  colorlinks=true,
  linkcolor=blue!70!black,
  urlcolor=cyan!70!black,
  pdftitle={ChemoInteract Part 1: Overview and Architecture},
  pdfauthor={Ronit},
  pdfstartview=FitH,
}

% ── Colours ──────────────────────────────────────────────────
\definecolor{codegreen}{rgb}{0,0.6,0}
\definecolor{codegray}{rgb}{0.5,0.5,0.5}
\definecolor{codepurple}{rgb}{0.58,0,0.82}
\definecolor{backcolour}{rgb}{0.97,0.97,0.97}
\definecolor{deepblue}{RGB}{0,55,120}
\definecolor{chemopurple}{RGB}{91,50,160}
\definecolor{chemogreen}{RGB}{16,140,100}
\definecolor{chemoorange}{RGB}{230,100,20}
\definecolor{noteyellow}{RGB}{255,248,200}
\definecolor{tipegreen}{RGB}{210,245,215}
\definecolor{warnorange}{RGB}{255,235,210}

% ── Listings ─────────────────────────────────────────────────
\lstdefinestyle{pythonstyle}{
  backgroundcolor=\color{backcolour},
  commentstyle=\color{codegreen}\itshape,
  keywordstyle=\color{chemopurple}\bfseries,
  numberstyle=\tiny\color{codegray},
  stringstyle=\color{chemogreen},
  basicstyle=\ttfamily\footnotesize,
  breaklines=true,
  captionpos=b,
  keepspaces=true,
  numbers=left,
  numbersep=5pt,
  showstringspaces=false,
  tabsize=4,
  frame=single,
  rulecolor=\color{gray!40},
  language=Python,
}
\lstset{style=pythonstyle}

% ── tcolorbox styles ─────────────────────────────────────────
\tcbset{
  notebox/.style={enhanced,breakable,colback=noteyellow,
    colframe=yellow!60!orange,fonttitle=\bfseries,title={Note},
    left=6pt,right=6pt,top=4pt,bottom=4pt},
  tipbox/.style={enhanced,breakable,colback=tipegreen,
    colframe=green!60!black,fonttitle=\bfseries,title={Tip},
    left=6pt,right=6pt},
  warnbox/.style={enhanced,breakable,colback=warnorange,
    colframe=orange!80!red,fonttitle=\bfseries,title={Warning},
    left=6pt,right=6pt},
}

% ── Header / Footer ──────────────────────────────────────────
\pagestyle{fancy}
\fancyhf{}
\fancyhead[L]{\small\textcolor{chemopurple}{\textbf{ChemoInteract --- Part 1}}}
\fancyhead[R]{\small\textcolor{gray}{Overview, Architecture \& Datasets}}
\fancyfoot[C]{\thepage}
\renewcommand{\headrulewidth}{0.5pt}

% ── Section styles ───────────────────────────────────────────
\titleformat{\section}{\Large\bfseries\color{deepblue}}{\thesection}{1em}{}
\titleformat{\subsection}{\large\bfseries\color{chemopurple}}{\thesubsection}{1em}{}
\titleformat{\subsubsection}{\normalsize\bfseries\color{chemogreen}}{\thesubsubsection}{1em}{}

\setstretch{1.2}
\setlength{\parskip}{6pt}

\begin{document}

% ── Title Block ──────────────────────────────────────────────
\begin{center}
{\Huge\bfseries\color{deepblue} ChemoInteract}\\[0.3cm]
{\Large\bfseries\color{chemopurple} Complete Documentation --- Part 1 of 3}\\[0.2cm]
{\large\textbf{System Overview, Directory Structure, File Reference, Workflow \& Datasets}}\\[0.4cm]
\begin{tcolorbox}[enhanced,width=14cm,colback=deepblue!5,colframe=deepblue,boxrule=1pt,arc=5pt]
\centering
\begin{tabular}{llll}
\textbf{Author:} & Ronit & \textbf{Version:} & 0.1.0 \\
\textbf{Framework:} & PyTorch 2.0+ & \textbf{Focus:} & System Architecture \\
\end{tabular}
\end{tcolorbox}
\end{center}

\tableofcontents
\newpage

% ================================================================
\section{Project Overview}\label{sec:overview}

\subsection{What is ChemoInteract?}
\textbf{ChemoInteract} is a modular deep learning library built on PyTorch. It provides a clean, standardized software framework for predicting what happens when two pharmaceutical drugs are administered together. It is a research-grade Python library --- not a web application or mobile app.

When a patient takes multiple medications simultaneously, those drugs can interact in complex biological ways. Some interactions cause adverse side effects; others improve therapeutic efficacy. ChemoInteract uses AI to \textit{predict} these interactions computationally \textit{before} conducting expensive wet-lab experiments.

\begin{tcolorbox}[notebox]
ChemoInteract is structured as a \textbf{Python package} (like NumPy or scikit-learn). Users install it, import modules, and invoke functions like \texttt{pipeline()} to train and evaluate models.
\end{tcolorbox}

\subsection{Core Pharmacology Problems Solved}

\begin{longtable}{>{\bfseries}p{4.5cm}p{5cm}p{5cm}}
\toprule
\textbf{Problem} & \textbf{Definition} & \textbf{Clinical Impact} \\
\midrule\endhead
Drug-Drug Interaction (DDI) &
Predicting if two drugs interact when co-administered &
Prevents adverse drug events in polypharmacy patients. \\
\midrule
Polypharmacy Side Effects &
Predicting specific side-effects resulting from drug pairs &
Identifies non-obvious adverse reactions not seen in single-drug trials. \\
\midrule
Drug Synergy &
Predicting if two drugs work \textit{better together} than independently &
Critical in oncology to design effective combination chemotherapy regimens. \\
\bottomrule
\end{longtable}

\subsection{Key Features}
\begin{itemize}
  \item \textbf{11 Deep Learning Models:} FC networks, GNNs, and Attention models in one repository.
  \item \textbf{5 Benchmark Datasets:} Auto-downloaded, parsed, and cached on disk.
  \item \textbf{One-Line Training:} Unified \texttt{pipeline()} replaces hundreds of lines of boilerplate.
  \item \textbf{Modular Design:} Swappable models, datasets, optimizers, and loss functions.
  \item \textbf{Automatic Hardware Acceleration:} CUDA GPU detection with seamless CPU fallback.
\end{itemize}

% ================================================================
\section{Folder and Package Structure}\label{sec:folders}

\subsection{Directory Tree}

\begin{lstlisting}[language=bash]
ChemoInteract/
|-- .git/                <- Git repository version control
|-- README.md            <- Overview and quickstart documentation
`-- ChemoInteract_Main/  <- Source code directory
    |-- ChemoInteract/   <- Main Python package
    |   |-- __init__.py        <- Package entry point
    |   |-- version.py         <- Version string ("0.1.0")
    |   |-- constants.py       <- Shared constants (e.g. node features)
    |   |-- compat.py          <- TorchDrug/RDKit compatibility patches
    |   |-- utils.py           <- Shared tensor & device utilities
    |   |-- loss.py            <- Custom loss functions (CASTER)
    |   |-- pipeline.py        <- Main pipeline() executor & Result dataclass
    |   |-- data/              <- Data loading & batching sub-package
    |   `-- models/            <- 11 deep learning model implementations
    |-- data_cleaning/   <- Preprocessing & cleaning scripts
    |-- dataset/         <- Pre-processed data JSON/CSV files
    |-- examples/        <- Demo scripts showing library usage
    `-- tests/           <- Unit test suite (pytest)
\end{lstlisting}

\subsection{Package Sub-Modules}

\begin{itemize}
  \item \textbf{\texttt{ChemoInteract.data}:} Contains dataset loaders (\texttt{DatasetLoader}, \texttt{DrugCombDB}, \texttt{TwoSides}, etc.), mini-batch iterator (\texttt{BatchGenerator}), feature containers (\texttt{DrugFeatureSet}, \texttt{ContextFeatureSet}), and triples wrapper (\texttt{LabeledTriples}).
  \item \textbf{\texttt{ChemoInteract.models}:} Contains abstract \texttt{Model} base class and all 11 model implementations (DeepSynergy, DeepDDI, MatchMaker, DeepDDS, DeepDrug, CASTER, GCNBMP, EPGCN-DS, MR-GNN, MHCADDI, SSI-DDI).
\end{itemize}

% ================================================================
\section{File-by-File Technical Reference}\label{sec:files}

\subsection{Core Package Files}

\begin{longtable}{>{\bfseries}p{4cm}p{10.5cm}}
\toprule
File & Purpose \& Functionality \\
\midrule\endhead
\texttt{version.py} & Defines \texttt{\_\_version\_\_ = "0.1.0"}. Tracks versioning in experiment results. \\
\midrule
\texttt{constants.py} & Dynamically calculates \texttt{TORCHDRUG\_NODE\_FEATURES} (69 atom features) via dummy molecule initialization. \\
\midrule
\texttt{compat.py} & Compatibility layer: patches missing RDKit \texttt{mplCanvas}, adds \texttt{.to(device)} method to TorchDrug \texttt{PackedGraph}, unifies atom feature dictionary keys across versions. \\
\midrule
\texttt{utils.py} & Provides \texttt{resolve\_device()}, \texttt{segment\_max()}, \texttt{segment\_sum()}, and numerically stable \texttt{segment\_softmax()}. \\
\midrule
\texttt{pipeline.py} & Houses \texttt{Result} dataclass and \texttt{pipeline()} function. Orchestrates device resolution, data splitting, model training, loss computation, evaluation, and logging. \\
\midrule
\texttt{loss.py} & Implements \texttt{CASTERSupervisedLoss}: a 3-part loss combining binary cross-entropy prediction loss, autoencoder reconstruction loss, and sparse dictionary regularisation. \\
\bottomrule
\end{longtable}

\begin{tcolorbox}[warnbox]
\textbf{Impact of removing \texttt{compat.py}:} TorchDrug \texttt{PackedGraph} objects will fail during GPU tensor transfer (\texttt{AttributeError: 'PackedGraph' object has no attribute 'to'}).
\end{tcolorbox}

% ================================================================
\section{Technology Stack}\label{sec:techstack}

\begin{longtable}{>{\bfseries}p{3.5cm}p{2.5cm}p{8.5cm}}
\toprule
\textbf{Library/Tool} & \textbf{Version} & \textbf{Role in Project} \\
\midrule\endhead
Python & 3.8+ & Primary programming language \\
PyTorch & 2.0+ & Deep learning tensors, autograd, GPU acceleration \\
TorchDrug & Latest & Molecular graphs (\texttt{PackedGraph}), GCN layers, readouts \\
RDKit & Latest & Cheminformatics, SMILES parsing, 256-bit Morgan fingerprints \\
scikit-learn & 1.0+ & Metrics (ROC-AUC, MSE, MAE), train/test splitting \\
Pandas & 1.3+ & DataFrames for triples and prediction storage \\
NumPy & 1.21+ & High-performance array operations \\
class-resolver & Latest & Dynamic string-to-class resolution for models \& loaders \\
pystow & Latest & Cross-platform directory cache management \\
tqdm & 4.0+ & Progress bar visualization during training \\
tabulate & 0.8+ & Console output table formatting \\
\bottomrule
\end{longtable}

% ================================================================
\section{Project Workflow and Execution Flow}\label{sec:workflow}

\subsection{End-to-End Execution Sequence}

\begin{enumerate}
  \item \textbf{Initialization:} User calls \texttt{pipeline(dataset="drugcombdb", model="deepsynergy", epochs=100)}.
  \item \textbf{Device Resolution:} \texttt{resolve\_device()} checks CUDA availability; assigns \texttt{torch.device("cuda")} or \texttt{torch.device("cpu")}.
  \item \textbf{Dataset Loading:} Resolver instantiates loader. Data files (\texttt{drug\_set.json}, \texttt{context\_set.json}, \texttt{labeled\_triples.csv}) are fetched and cached via \texttt{@lru\_cache}.
  \item \textbf{Data Splitting:} \texttt{LabeledTriples} splits 80\% train, 20\% test (\texttt{random\_state=42}).
  \item \textbf{Batching:} Two \texttt{BatchGenerator} objects instantiated for training and testing loops.
  \item \textbf{Model \& Optimizer Setup:} Model instantiated and moved to target device via \texttt{model.to(device)}. Adam optimizer initialized ($\alpha=0.001$).
  \item \textbf{Training Loop:}
  \begin{enumerate}[label=\roman*.]
    \item \texttt{BatchGenerator} yields mini-batch \texttt{DrugPairBatch}.
    \item \texttt{optimizer.zero\_grad()} resets gradients.
    \item \texttt{model.unpack(batch)} extracts required input tensors.
    \item \texttt{model.forward(...)} computes predictions $\hat{y} \in [0, 1]$.
    \item \texttt{loss\_fn(pred, target)} calculates scalar loss.
    \item \texttt{loss.backward()} computes autograd gradients.
    \item \texttt{optimizer.step()} updates model parameters.
  \end{enumerate}
  \item \textbf{Evaluation:} Model set to \texttt{model.eval()}. Predictions generated on test set inside \texttt{torch.no\_grad()}. ROC-AUC score computed.
  \item \textbf{Output:} Returns populated \texttt{Result} object with training history and metrics.
\end{enumerate}

% ================================================================
\section{Datasets and Feature Representations}\label{sec:datasets}

\subsection{Supported Benchmark Datasets}

\begin{longtable}{>{\bfseries}p{3.5cm}p{2.5cm}p{2.5cm}p{2.5cm}p{3cm}}
\toprule
\textbf{Dataset} & \textbf{Task} & \textbf{Loader Type} & \textbf{Context Type} & \textbf{Primary Source} \\
\midrule\endhead
DrugCombDB & Synergy & Remote & Cell lines & Anti-cancer screen \\
DrugComb & Synergy & Remote & Cell lines & FIMM database \\
TwoSides & Side-Effects & Remote & Side effect ID & Polypharmacy reporting \\
DrugbankDDI & DDI & Remote & Interaction type & FDA DrugBank \\
OncoPolyPharmacology & Synergy & Local (TDC) & Cancer cell lines & O'Neil 2016 screen \\
\bottomrule
\end{longtable}

\subsection{Feature Storage Format}
Every dataset directory maintains three standard files:
\begin{lstlisting}[language=bash]
dataset/<dataset_name>/
|-- drug_set.json        <- {"drug_id": {"smiles": "...", "features": [256 floats]}}
|-- context_set.json     <- {"context_id": [float vector]}
`-- labeled_triples.csv  <- columns: drug_1, drug_2, context, label
\end{lstlisting}

\subsection{Negative Sampling Strategy}
Database records only provide positive (known) interactions. Negative samples (label = 0.0) are generated via rejection sampling under the \textbf{Closed-World Assumption (CWA)}: any unrecorded drug-drug-context pair is assumed non-interacting.

\end{document}
